%% figures.tex: example plots for the template. ALL DATA IS INVENTED; replace it with yours.
%% Each plot is a macro taking the panel width and height, so poster.tex can resize it.
%% Styles come from posterplots.tex.

%% ---------------------------------------------------------------
%% Data (invented). grid{S,L}{A,B,C}: model size x dataset; row{A,B,C}: baselines per dataset.
%% ---------------------------------------------------------------
\pgfplotstableread{
step ours variant a oracle
20 23.2 23.5 22.5 24.5
40 33.9 34.0 34.7 36.8
60 49.0 47.7 47.2 51.0
80 58.1 59.3 58.2 63.0
100 64.8 63.7 63.5 69.6
120 69.7 68.5 67.5 73.0
140 71.2 70.3 56.8 75.9
160 72.0 71.7 11.1 75.1
180 71.5 70.7 1.1 76.0
200 71.6 70.0 0.0 76.7
220 71.3 70.6 0.4 75.1
240 72.0 60.2 0.0 75.8
260 71.9 11.2 0.4 76.0
280 70.8 1.3 0.6 76.6
300 73.2 0.3 0.1 75.2
}\gridSA

\pgfplotstableread{
step ours variant a oracle
20 18.4 17.0 16.8 18.7
40 25.7 25.5 26.7 27.6
60 35.6 36.3 36.8 40.5
80 43.9 42.3 44.0 49.1
100 49.6 50.5 49.7 55.2
120 53.3 52.5 52.6 58.3
140 54.6 53.7 50.4 59.9
160 55.4 54.2 42.1 59.3
180 55.5 52.2 8.7 60.5
200 53.9 55.0 0.9 59.9
220 55.2 45.5 0.1 60.7
240 54.5 8.7 0.9 60.0
260 54.3 1.3 1.3 59.7
280 53.9 0.0 0.0 59.8
300 56.1 0.0 1.1 59.2
}\gridSB

\pgfplotstableread{
step ours variant a oracle
20 9.4 10.3 10.4 11.2
40 15.2 14.6 14.1 16.2
60 20.6 20.3 19.9 22.1
80 26.2 25.3 25.7 27.4
100 28.1 27.2 26.6 30.9
120 29.6 29.3 29.6 30.3
140 29.6 29.4 28.9 32.7
160 30.5 29.7 29.1 32.5
180 32.9 29.2 28.2 32.9
200 30.8 28.5 21.6 32.7
220 31.8 27.1 4.8 33.6
240 31.7 29.2 0.0 32.8
260 30.7 28.0 1.0 31.4
280 31.9 25.8 0.6 32.1
300 31.1 27.8 0.0 33.1
}\gridSC

\pgfplotstableread{
step ours variant a oracle
20 27.0 26.1 25.6 28.5
40 39.8 38.2 40.5 40.2
60 54.9 53.3 53.0 57.3
80 67.0 66.6 63.8 69.2
100 74.9 72.6 71.7 77.2
120 79.2 77.9 77.6 81.5
140 79.8 76.9 78.5 84.7
160 79.8 79.0 79.4 84.4
180 79.2 78.2 78.7 84.4
200 81.2 76.5 80.3 84.3
220 81.9 76.5 80.1 84.9
240 80.4 75.1 76.0 85.1
260 82.1 73.0 37.4 84.8
280 79.5 73.1 3.4 84.6
300 81.0 72.1 0.2 85.8
}\gridLA

\pgfplotstableread{
step ours variant a oracle
20 20.1 20.8 20.8 22.1
40 29.3 30.1 27.2 30.6
60 39.9 41.8 39.0 43.5
80 51.0 50.3 49.0 53.8
100 58.4 56.1 55.6 60.3
120 59.7 57.7 57.4 63.4
140 59.8 58.9 60.0 64.5
160 61.6 59.9 59.8 63.9
180 60.6 58.1 60.7 64.5
200 61.2 57.3 58.6 64.9
220 61.0 57.6 57.8 65.2
240 61.5 54.9 60.5 64.8
260 60.2 55.1 57.9 64.7
280 62.6 55.8 30.6 65.2
300 63.1 55.0 2.8 63.7
}\gridLB

\pgfplotstableread{
step ours variant a oracle
20 13.7 13.9 12.1 13.0
40 20.8 17.2 18.7 21.2
60 26.0 26.7 27.1 27.4
80 33.4 33.0 29.8 33.8
100 37.0 36.6 33.2 37.6
120 37.8 37.3 34.8 39.6
140 38.7 37.2 36.3 41.1
160 40.3 35.5 33.8 41.0
180 41.2 37.9 32.4 41.2
200 38.4 38.0 31.9 40.5
220 41.0 38.3 29.4 40.6
240 40.2 36.4 29.4 40.4
260 41.7 36.7 27.8 40.2
280 41.4 36.1 29.6 41.5
300 39.3 35.0 25.1 40.6
}\gridLC

\pgfplotstableread{
step a b c
20 23.4 23.2 16.0
40 34.5 34.0 19.6
60 48.7 47.0 22.6
80 60.0 57.8 26.6
100 65.6 64.3 32.3
120 66.9 67.6 38.1
140 35.4 68.0 43.8
160 3.6 69.9 48.5
180 0.5 69.0 51.5
200 0.1 68.6 57.2
220 0.0 56.1 58.5
240 1.3 11.5 60.2
260 0.0 1.0 63.2
280 0.1 1.2 64.4
300 0.0 0.5 65.8
}\rowA

\pgfplotstableread{
step a b c
20 18.7 18.0 11.4
40 26.3 26.1 14.1
60 37.6 35.9 17.7
80 45.1 45.8 21.3
100 48.4 48.8 26.1
120 26.3 51.8 29.0
140 2.1 51.3 32.4
160 0.4 53.5 36.7
180 0.0 51.5 39.8
200 0.2 26.5 41.8
220 0.5 1.2 43.5
240 0.0 0.0 45.2
260 0.0 0.0 47.1
280 0.5 0.0 47.2
300 0.0 1.5 48.3
}\rowB

\pgfplotstableread{
step a b c
20 10.9 8.7 6.1
40 13.5 14.6 8.1
60 19.2 19.4 9.2
80 24.2 22.5 9.5
100 28.0 24.5 12.3
120 30.1 26.3 14.9
140 30.6 26.1 15.4
160 15.0 23.4 17.6
180 1.1 23.1 20.5
200 0.0 21.0 21.4
220 0.0 20.6 22.4
240 0.0 20.3 23.5
260 0.0 18.0 23.6
280 0.0 16.6 24.2
300 0.0 16.4 24.6
}\rowC


%% ---------------------------------------------------------------
%% Learning curves, one row of small multiples (block 3).
%% Markers on every second step: the guide allows markers only with few x positions.
%% ---------------------------------------------------------------
\newcommand{\RowLines}[1]{%
  \addplot[s c, mark repeat=2, mark phase=1] table[x=step, y=c] {#1};
  \addplot[s b, mark repeat=2, mark phase=2] table[x=step, y=b] {#1};
  \addplot[s a, mark repeat=2, mark phase=1] table[x=step, y=a] {#1};
}
\newcommand{\CurvesRow}[2]{%
\begin{tikzpicture}
\begin{groupplot}[
  group style={group size=3 by 1, horizontal sep=2.0cm, ylabels at=edge left},
  posteraxis, width=#1, height=#2,
  xmin=10, xmax=310, xtick={100,200,300}, ymin=0,
  ylabel={accuracy (\%)},
]
\nextgroupplot[title={Dataset A}, ymax=90, ytick={0,20,40,60,80}]                         \RowLines{\rowA}
\nextgroupplot[title={Dataset B}, ymax=70, ytick={0,20,40,60}, xlabel={training step}]  \RowLines{\rowB}
\nextgroupplot[title={Dataset C}, ymax=45, ytick={0,10,20,30,40}]                       \RowLines{\rowC}
\end{groupplot}
\end{tikzpicture}}

%% ---------------------------------------------------------------
%% Learning curves, 2 x 3 small multiples (block 6). Ours is drawn last, on top.
%% One y-label per row, one shared x-label, row labels on the left.
%% ---------------------------------------------------------------
\newcommand{\GridLines}[1]{%
  \addplot[s oracle]                            table[x=step, y=oracle]  {#1};
  \addplot[s a, mark repeat=2, mark phase=2]    table[x=step, y=a]       {#1};
  \addplot[s variant]                           table[x=step, y=variant] {#1};
  \addplot[s ours, mark repeat=2, mark phase=1] table[x=step, y=ours]    {#1};
}
\newcommand{\CurvesGrid}[2]{%
\begin{tikzpicture}
\begin{groupplot}[
  group style={group size=3 by 2, horizontal sep=2.3cm, vertical sep=1.4cm,
               ylabels at=edge left, xticklabels at=edge bottom},
  posteraxis, width=#1, height=#2,
  xmin=10, xmax=310, xtick={100,200,300}, ymin=0,
  ylabel={accuracy (\%)},
]
\nextgroupplot[title={Dataset A}, ymax=90, ytick={0,20,40,60,80}]        \GridLines{\gridSA}
\nextgroupplot[title={Dataset B}, ymax=70, ytick={0,20,40,60}]           \GridLines{\gridSB}
\nextgroupplot[title={Dataset C}, ymax=45, ytick={0,10,20,30,40}]        \GridLines{\gridSC}
\nextgroupplot[ymax=90, ytick={0,20,40,60,80}]                           \GridLines{\gridLA}
\nextgroupplot[ymax=70, ytick={0,20,40,60}, xlabel={training step}]      \GridLines{\gridLB}
\nextgroupplot[ymax=45, ytick={0,10,20,30,40}]                           \GridLines{\gridLC}
\end{groupplot}
\node[rotate=90, anchor=south, font=\sffamily\figlabel, inkc] at ($(group c1r1.west)+(-3.0cm,0)$) {Model-S};
\node[rotate=90, anchor=south, font=\sffamily\figlabel, inkc] at ($(group c1r2.west)+(-3.0cm,0)$) {Model-L};
\end{tikzpicture}}

%% ---------------------------------------------------------------
%% Grouped bars from zero (block 8). Values printed for ours only (with 4 bars per group, a second
%% set of labels collides; with 2-3 bars, label the strongest baseline too).
%% Two-level tick labels (size, then family) keep long category names from colliding.
%% ---------------------------------------------------------------
\newcommand{\BarsPlot}[2]{%
\begin{tikzpicture}
\begin{groupplot}[
  group style={group size=2 by 1, horizontal sep=1.6cm, ylabels at=edge left, yticklabels at=edge left},
  posterbars, width=#1, height=#2,
  ybar=0.14cm, /pgf/bar width=0.68cm,
  symbolic x coords={s,m,l,y}, xtick=data, xticklabels={S, M, L, S},
  enlarge x limits=0.16,
  ymin=0, ymax=100, ytick={0,20,40,60,80},
  ylabel={accuracy (\%)},
  clip=false,   % lets the family labels sit under the ticks; bars stay inside the limits
  valuelabels/.style={nodes near coords={\pgfmathprintnumber[fixed,precision=0]{\pgfplotspointmeta}}},
]
\nextgroupplot[title={Dataset A}]
\famlabel{m}{Model X}\famlabel{y}{Model Y}
\addplot[b ours, valuelabels] coordinates {(s,71.2) (m,76.5) (l,81.0) (y,69.8)};
\addplot[b a]                 coordinates {(s,58.4) (m,63.1) (l,70.2) (y,61.5)};
\addplot[b b]                 coordinates {(s,52.0) (m,57.7) (l,64.9) (y,55.3)};
\addplot[b oracle]            coordinates {(s,76.0) (m,80.3) (l,84.1) (y,74.6)};
\nextgroupplot[title={Dataset B}]
\famlabel{m}{Model X}\famlabel{y}{Model Y}
\addplot[b ours, valuelabels] coordinates {(s,52.6) (m,58.0) (l,63.4) (y,50.1)};
\addplot[b a]                 coordinates {(s,41.9) (m,47.5) (l,55.0) (y,43.2)};
\addplot[b b]                 coordinates {(s,37.4) (m,42.8) (l,49.6) (y,39.0)};
\addplot[b oracle]            coordinates {(s,57.3) (m,61.9) (l,66.0) (y,55.4)};
\end{groupplot}
\end{tikzpicture}}

%% ---------------------------------------------------------------
%% Ablation as a dumbbell (block 9): the claim is a before/after change per variant.
%% Ablations of ours use tints of the ours hue (guide, section 4). If an ablation equals another
%% method shown on the poster, give it that method's colour instead.
%% ---------------------------------------------------------------
\newcommand{\AblationPlot}[2]{%
\begin{tikzpicture}
\pgfplotsset{abl/.style={
  posteraxis, width=#1, height=#2,
  y dir=reverse, ymin=0, ymax=3, ytick={0,1,2,3}, enlarge y limits={abs=0.6},
  ymajorgrids=false, clip=false,
}}
\begin{axis}[abl, name=left,
  xmin=-4, xmax=84, xtick={0,20,40,60,80},   % dots at 0 sit inside the frame; ticks stay in range
  yticklabels={{Full method (ours)}, {w/o component 1}, {w/o component 2}, {w/o component 3}},
  xlabel={Dataset A accuracy (\%)},
]
\Dumbbell{oursc}{0}{74.0}{71.2}
\Dumbbell{oursc!70!white}{1}{73.1}{48.6}
\Dumbbell{oursc!55!white}{2}{72.5}{30.2}
\Dumbbell{oursc!40!white}{3}{70.8}{12.4}
\node[anchor=east, font=\sffamily\figvalue, inkc, xshift=-0.55cm] at (axis cs:71.2,0) {71.2};
\end{axis}
\begin{axis}[abl, at={($(left.east)+(2.0cm,0)$)}, anchor=west,
  yticklabels={}, xlabel={Dataset B accuracy (\%)}, xmin=-3, xmax=63, xtick={0,20,40,60},
]
\Dumbbell{oursc}{0}{55.0}{52.6}
\Dumbbell{oursc!70!white}{1}{54.2}{35.1}
\Dumbbell{oursc!55!white}{2}{53.8}{20.4}
\Dumbbell{oursc!40!white}{3}{52.9}{8.3}
\node[anchor=east, font=\sffamily\figvalue, inkc, xshift=-0.55cm] at (axis cs:52.6,0) {52.6};
\end{axis}
\end{tikzpicture}}
